\setcounter{section}{6}

  \zwischenueberschrift{Kovariante Ableitung}

Zu einer differenzierbaren Hyperfläche

\mavergleichskette
{\vergleichskette
{ Y
}
{ = }{ \R^n
}
{  }{
}
{    }{
}
{   }{
}
}
{}{}{}
möchte man Ableitungskonzepte, die im umgebenden Raum $\R^n$ einfach definiert sind, mit einem natürlichen Bezug zu $Y$ definieren. Man spricht von kovarianter Ableitung, das wichtigste Hilfsmittel ist die orthogonale Projektion
\maabbdisp {\pi_P} {\R^n } { T_PY
} {}
zu einem Punkt

\mavergleichskette
{\vergleichskette
{ P
}
{ \in }{ Y
}
{  }{
}
{    }{
}
{   }{
}
}
{}{}{.}

\inputdefinition
{}
{

Es sei

\mavergleichskette
{\vergleichskette
{ Y
}
{ \subseteq }{ W
}
{ \subseteq }{ \R^n
}
{    }{
}
{   }{
}
}
{}{}{,}
$W$
\definitionsverweis {offen}{}{,}
eine
\definitionsverweis {differenzierbare Hyperfläche}{}{,}
sei

\mavergleichskette
{\vergleichskette
{ P
}
{ \in }{ Y
}
{  }{
}
{    }{
}
{   }{
}
}
{}{}{}
ein Punkt und

\mavergleichskette
{\vergleichskette
{ v
}
{ \in }{ T_PY
}
{  }{
}
{    }{
}
{   }{
}
}
{}{}{}
ein
\definitionsverweis {Tangentialvektor}{}{.}
Es sei
\maabbdisp {F} { U} { \R^n
} {}
ein differenzierbares Vektorfeld, das auf einer offenen Umgebung

\mavergleichskette
{\vergleichskette
{ P
}
{ \in }{ U
}
{ \subseteq }{ \R^n
}
{    }{
}
{   }{
}
}
{}{}{}
definiert sei und das auf $Y \cap U$  tangential an $Y$ sei. Dann nennt man

\mavergleichskettedisp
{\vergleichskette
{ \nabla_v F
}
{ \defeq} { \pi_P   \left(  \left(  D F   \right)_{ P } \left(  v  \right)  \right)
}
{ } {
}
{ } {
}
{ } {
}
}
{}{}{,}
wobei
\maabbdisp {\pi_P} {\R^n } { T_PY
} {}
die
\definitionsverweis {orthogonale Projektion}{}{}
bezeichnet, die
\definitionswort {kovariante Ableitung}{}
von $F$ in Richtung $v$.

}

Es liegt insgesamt die Situation

\mathdisp {T_PY \longrightarrow \R^n \stackrel{  \left(D F \right)_{ P} }{\longrightarrow} \R^n\stackrel{ \pi_P }{\longrightarrow} T_PY} {  }

vor und $\nabla_v F$ ist das Ergebnis, wenn man vorne den Tangentialvektor

\mavergleichskette
{\vergleichskette
{ v
}
{ \in }{ T_PY
}
{  }{
}
{    }{
}
{   }{
}
}
{}{}{}
einsetzt. Wenn $N$ ein Einheitsnormalenfeld bezeichnet, so ist die kovariante Ableitung gleich

\mavergleichskettedisp
{\vergleichskette
{ \nabla_v F
}
{ =} { \left(  D F   \right)_{ P } \left(  v  \right) -  \left\langle  \left(  D F   \right)_{ P } \left(  v  \right)  ,  N(P)  \right\rangle N(P)
}
{ } {
}
{ } {
}
{ } {
}
}
{}{}{,}
da man ja so das orthogonale Komplement berechnet.

\inputdefinition
{}
{

Es sei

\mavergleichskette
{\vergleichskette
{ Y
}
{ \subseteq }{ W
}
{ \subseteq }{ \R^n
}
{    }{
}
{   }{
}
}
{}{}{,}
$W$
\definitionsverweis {offen}{}{,}
eine
\definitionsverweis {differenzierbare Hyperfläche}{}{}
und sei
\maabbdisp {G} { Y } { TY
} {}
ein tangentiales Vektorfeld auf $Y$. Es sei
\maabbdisp {F} { U} { \R^n
} {}
ein differenzierbares Vektorfeld, das auf einer offenen Umgebung

\mavergleichskette
{\vergleichskette
{ P
}
{ \in }{ U
}
{ \subseteq }{ \R^n
}
{    }{
}
{   }{
}
}
{}{}{}
definiert sei und das auf $Y \cap U$  tangential an $Y$ sei. Dann nennt man das tangentiale Vektorfeld
\maabbdisp {\nabla_G F} { Y \cap U } { T(Y \cap U)
} {,}
das durch

\mavergleichskettedisp
{\vergleichskette
{ \left(  \nabla_G F  \right) (P)
}
{ \defeq} { \nabla_{G(P)} F }
{ } {
}
{ } {
}
{ } {
}
}
{}{}{}
definiert ist, die
\definitionswort {kovariante Ableitung}{}
von $F$ in Richtung $G$.

}

__NOEDITSECTION__

\inputfaktbeweis
{Hyperfläche/Tangentiales Vektorfeld/Bezüglich Tangentenvektor/Kovariante Ableitung/Eigenschaften/Fakt}
{Lemma}
{}
{

\faktsituation {Es sei

\mavergleichskette
{\vergleichskette
{ Y
}
{ \subseteq }{ W
}
{ \subseteq }{\R^n
}
{    }{
}
{   }{
}
}
{}{}{,}
$W$
\definitionsverweis {offen}{}{,}
eine
\definitionsverweis {differenzierbare Hyperfläche}{}{,}
sei

\mavergleichskette
{\vergleichskette
{ P
}
{ \in }{ Y
}
{  }{
}
{    }{
}
{   }{
}
}
{}{}{}
ein Punkt.}
\faktuebergang {Dann gelten folgende Aussagen.}
\faktfolgerung {\aufzaehlungdrei{Zu einem fixierten differenzierbaren
\definitionsverweis {tangentialen Vektorfeld}{}{}
$F$ ist die Zuordnung
\maabbeledisp {} {T_PY} { T_PY
} { v } { \nabla_v F
} {,}
\definitionsverweis {linear}{}{.}
}{Zu einem fixierten Tangentialvektor

\mavergleichskette
{\vergleichskette
{ v
}
{ \in }{ T_PY
}
{  }{
}
{    }{
}
{   }{
}
}
{}{}{}
ist die Zuordnung
\maabbdisp {} { F } { \nabla_v F
} {,}
die einem differenzierbaren Vektorfeld die kovariante Ableitung längs  $v$ zuordnet, linear.
}{Zu einer differenzierbaren Funktion
\maabb {g} { U } {\R
} {}
und einem differenzierbaren tangentialen Vektorfeld $F$ ist

\mavergleichskettedisp
{\vergleichskette
{ \nabla_v (gF)
}
{ =} { g (P) ( \nabla_v F) (P) + \left(  D_{v}  g   \right) \left(   P   \right) F(P)
}
{ } {
}
{ } {
}
{ } {
}
}
{}{}{.}
}}
\faktzusatz {}
\faktzusatz {}

}
{

\aufzaehlungdrei{Dies folgt unmittelbar aus der Beschreibung

\mavergleichskettedisp
{\vergleichskette
{ \nabla_v F
}
{ =} { \left(  D F   \right)_{ P } \left(  v  \right) - \left\langle  \left(  D F   \right)_{ P } \left(  v  \right)  ,  N(P)  \right\rangle N(P)
}
{ } {
}
{ } {
}
{ } {
}
}
{}{}{,}
da die beiden Summanden linear in $v$ sind.
}{Dies folgt auch aus der Beschreibung

\mavergleichskettedisp
{\vergleichskette
{ \nabla_v F
}
{ =} { \left(  D F   \right)_{ P } \left(  v  \right) - \left\langle  \left(  D F   \right)_{ P } \left(  v  \right)  ,  N(P)  \right\rangle N(P)
}
{ } {
}
{ } {
}
{ } {
}
}
{}{}{,}
da bei fixiertem Vektor $v$ die beiden Summanden gemäß
[[Richtungsableitung/R/Elementare Eigenschaften/Fakt|Lemma 43.6 (Analysis (Osnabrück 2021-2023))]]
linear von $F$ abhängen.
}{Nach
[[Totale Differenzierbarkeit/K/Produktregel/Fakt|Lemma 45.10 (Analysis (Osnabrück 2021-2023))]]
ist

\mavergleichskettedisp
{\vergleichskette
{ \left(D( g F)\right)_{ P }
}
{ =} { g(P) \cdot \left(D F \right)_{ P } + \left(D g \right)_{ P } \cdot F (P)
}
{ } {
}
{ } {
}
{ } {
}
}
{}{}{.}
Daher ist

\mavergleichskettealignhandlinks
{\vergleichskettealignhandlinks
{ \nabla_v (gF)
}
{ =} { \left(  DgF  \right)_{ P } \left(  v  \right) - \left\langle  \left(  DgF  \right)_{ P } \left(  v  \right)  ,  N(P)  \right\rangle N(P)
}
{ =} { g(P) \cdot \left(  D F   \right)_{ P } \left(  v  \right) + \left(  D g   \right)_{ P } \left(  v  \right) \cdot F (P) - \left\langle  g(P) \cdot \left(  D F   \right)_{ P } \left(  v  \right) + \left(  D g   \right)_{ P } \left(  v  \right) \cdot F (P)  ,  N(P)  \right\rangle N(P)
}
{ =} { g(P) \cdot \left(  D F   \right)_{ P } \left(  v  \right) - \left\langle  g(P) \cdot \left(  D F   \right)_{ P } \left(  v  \right)  ,  N(P)  \right\rangle N(P)        + \left(  D g   \right)_{ P } \left(  v  \right) \cdot F (P) - \left\langle  \left(  D g   \right)_{ P } \left(  v  \right) \cdot F (P)  ,  N(P)  \right\rangle N(P)
}
{ =} { g (P) ( \nabla_v F) (P) + \left(  D_{v}  g   \right) \left(   P   \right) \cdot F(P)
}
}
{}
{}{,}
da das Vektorfeld $F$ senkrecht auf $N$ steht.
}

}

__NOEDITSECTION__

\inputfaktbeweis
{Hyperfläche/Tangentiales Vektorfeld/Bezüglich tangentialem Vektorfeld/Kovariante Ableitung/Eigenschaften/Fakt}
{Lemma}
{}
{

\faktsituation {Es sei

\mavergleichskette
{\vergleichskette
{ Y
}
{ \subseteq }{ W
}
{ \subseteq }{\R^n
}
{    }{
}
{   }{
}
}
{}{}{,}
$W$
\definitionsverweis {offen}{}{,}
eine
\definitionsverweis {differenzierbare Hyperfläche}{}{.}}
\faktuebergang {Dann gelten für tangentiale Vektorfelder auf $Y$ folgende Aussagen.}
\faktfolgerung {\aufzaehlungdrei{Zu einem fixierten differenzierbaren tangentialen Vektorfeld $F$ ist die Zuordnung
\mathl{G \mapsto \nabla_G F}{}
\definitionsverweis {linear}{}{}
in $G$. Ferner ist für eine stetige Funktion
\maabb {g} { Y } {\R
} {}

\mavergleichskettedisp
{\vergleichskette
{ \nabla_{gG} F
}
{ =} { g \nabla_G F
}
{ } {
}
{ } {
}
{ } {
}
}
{}{}{.}
}{Zu einem fixierten tangentialen Vektorfeld $G$ ist die Zuordnung
\mathl{F \mapsto \nabla_G F}{,} die einem tangentialen differenzierbaren Vektorfeld die kovariante Ableitung längs $G$ zuordnet, linear.
}{Zu einer differenzierbaren Funktion
\maabb {g} { U } {\R
} {}
und einem differenzierbaren tangentialen Vektorfeld $F$ ist

\mavergleichskettedisp
{\vergleichskette
{ \nabla_G (gF)
}
{ =} { g (P) ( \nabla_G F) (P) + (\nabla_G g )(P) F(P)
}
{ } {
}
{ } {
}
{ } {
}
}
{}{}{.}
}}
\faktzusatz {}
\faktzusatz {}

}
{

Alle Aussagen folgen aus
[[Hyperfläche/Tangentiales Vektorfeld/Bezüglich Tangentenvektor/Kovariante Ableitung/Eigenschaften/Fakt|Lemma 6.3]].

}

In der letzten Gleichung bedeutet

\mavergleichskettedisp
{\vergleichskette
{ \left(  \nabla_G g  \right) (P)
}
{ =} { \left(  D_{ G(P)}  g   \right) \left(   P   \right)
}
{ } {
}
{ } {
}
{ } {
}
}
{}{}{,}
es wird also von der Funktion $g$ die Richtungsableitung in Richtung $G(P)$ genommen.

\inputdefinition
{}
{

Es sei

\mavergleichskette
{\vergleichskette
{ Y
}
{ \subseteq }{ W
}
{ \subseteq }{\R^n
}
{    }{
}
{   }{
}
}
{}{}{,} $W$
\definitionsverweis {offen}{}{,}
eine
\definitionsverweis {differenzierbare Hyperfläche}{}{,}
sei
\maabbdisp {\gamma} {[a,b]} {Y
} {}
eine
\definitionsverweis {differenzierbare Kurve}{}{}
und sei
\maabbdisp {F} { I } { \R^n
} {}
ein differenzierbares Vektorfeld längs $I$, das tangential an $Y$ sei. Dann nennt man

\mavergleichskettedisp
{\vergleichskette
{ (\nabla_\gamma F) (t)
}
{ \defeq} { \pi_{\gamma(t)} F' (t)
}
{ } {
}
{ } {
}
{ } {
}
}
{}{}{,}
die
\definitionswort {kovariante Ableitung}{}
von $F$ längs $\gamma$.

}

Man leitet also einfach das Vektorfeld $F$, das ja eine vektorwertige Kurve mit

\mavergleichskette
{\vergleichskette
{ F(t)
}
{ \in }{ T_{\gamma(t)}Y
}
{  }{
}
{    }{
}
{   }{
}
}
{}{}{}
ist, als Kurve im $\R^n$ ab und nimmt vom Ergebnis die orthogonale Projektion auf den Tangentialraum. Manchmal schreibt man auch $\nabla_{\gamma'} F$ oder $\nabla_{\partial} F$\zusatzfussnote {Zu
[[Hyperfläche/Tangentiales Vektorfeld/Bezüglich tangentialem Vektorfeld/Kovariante Ableitung/Definition|Definition 6.2]]
und
[[Hyperfläche/Tangentiales Vektorfeld/Längs einer Kurve/Kovariante Ableitung/Definition|Definition 6.5]]
gibt es die folgende Verallgemeinerung. Es sei $Y$ die differenzierbare Hyperfläche
\zusatzklammer {oder eine differenzierbare Mannigfaltigkeit mit einem Zusammenhang auf dem Tangentialbündel} {} {,}
sei $Z$ eine differenzierbare Mannigfaltigkeit, sei
\maabb {\varphi} { Z } { Y
} {}
eine differenzierbare Abbildung, sei
\maabbdisp {F} { Z } { TY = \biguplus_{P \in Y} T_PY
} {}
ein differenzierbares tangentiales Vektorfelder längs $\varphi$ und $G$ ein Vektorfeld auf $Z$. Dann ist
\zusatzklammer {zu

\mavergleichskette
{\vergleichskette
{ Q
}
{ \in }{ Z
}
{  }{
}
{    }{
}
{   }{
}
}
{}{}{}
und

\mavergleichskette
{\vergleichskette
{ P
}
{ = }{ \varphi(Q)
}
{  }{
}
{    }{
}
{   }{
}
}
{}{}{}} {} {}
die kovariante Ableitung von $F$ bezüglich $G$
\zusatzklammer {längs $\varphi$} {} {}
durch

\mavergleichskettedisp
{\vergleichskette
{ ( \nabla_G F )(Q)
}
{ =} { \pi_{ P  } \left(D F \right)_{ Q}  (G(Q))
}
{ } {
}
{ } {
}
{ } {
}
}
{}{}{}
definiert. $\nabla_GF$ ist wieder eine Abbildung
\maabb {} { Z } { TY
} {.}
In der ersten Definition ist

\mavergleichskette
{\vergleichskette
{ Z
}
{ = }{ Y
}
{  }{
}
{    }{
}
{   }{
}
}
{}{}{}
und $\varphi$ die Identität, in der zweiten Definition ist

\mavergleichskette
{\vergleichskette
{ Z
}
{ = }{ I
}
{  }{
}
{    }{
}
{   }{
}
}
{}{}{}
ein Intervall,

\mavergleichskette
{\vergleichskette
{ \varphi
}
{ = }{ \gamma
}
{  }{
}
{    }{
}
{   }{
}
}
{}{}{}
der Weg und

\mavergleichskette
{\vergleichskette
{ G
}
{ = }{ \partial
}
{  }{
}
{    }{
}
{   }{
}
}
{}{}{}
das konstante Vektorfeld auf dem Intervall $I$ mit der Richtung $1$.} {} {.}

  \zwischenueberschrift{Parallele Vektorfelder}

\inputdefinition
{}
{

Es sei

\mavergleichskette
{\vergleichskette
{ Y
}
{ \subseteq }{ W
}
{ \subseteq }{ \R^n
}
{    }{
}
{   }{
}
}
{}{}{,}
$W$
\definitionsverweis {offen}{}{,}
eine
\definitionsverweis {differenzierbare Hyperfläche}{}{}
und sei
\maabb {\gamma} { I } { Y
} {}
eine
\definitionsverweis {differenzierbare Kurve}{}{.}
Man sagt, dass ein längs $\gamma$ definiertes differenzierbares tangentiales Vektorfeld
\maabb {F} { I } {\R^n
} {}
\definitionswort {parallel längs}{}
$\gamma$ ist, wenn

\mavergleichskettedisp
{\vergleichskette
{ (\nabla_{\gamma} F)(t)
}
{ =} { 0
}
{ } {
}
{ } {
}
{ } {
}
}
{}{}{}
für alle

\mavergleichskette
{\vergleichskette
{ t
}
{ \in }{ I
}
{  }{
}
{    }{
}
{   }{
}
}
{}{}{}
gilt.

}

Ein Vektorfeld ist genau dann parallel längs $\gamma$, wenn
\mathl{F'(t)}{} stets senkrecht zum Tangentialraum
\mathl{T_{\gamma(t)}Y}{} ist, also zur Normalengerade gehört.

__NOEDITSECTION__
\inputfaktbeweis
{Hyperfläche/Kurve/Paralleles Vektorfeld/Eigenschaften/Fakt}
{Lemma}
{}
{

\faktsituation {Es sei

\mavergleichskette
{\vergleichskette
{ Y
}
{ \subseteq }{ W
}
{ \subseteq }{ \R^n
}
{    }{
}
{   }{
}
}
{}{}{,}
$W$
\definitionsverweis {offen}{}{,}
eine
\definitionsverweis {differenzierbare Hyperfläche}{}{}
und sei
\maabb {\gamma} { I } { Y
} {}
eine
\definitionsverweis {differenzierbare Kurve}{}{.}}
\faktuebergang {Dann gelten folgende Aussagen.}
\faktfolgerung {\aufzaehlungzwei {Zu einem längs $\gamma$
\definitionsverweis {parallelen Vektorfeld}{}{}
$F$ und

\mavergleichskette
{\vergleichskette
{ a
}
{ \in }{ \R
}
{  }{
}
{    }{
}
{   }{
}
}
{}{}{}
ist auch
\mathl{aF}{} ein paralleles Vektorfeld.
} {Zu längs $\gamma$ parallelen Vektorfelder
$F , G$ ist auch
\mathl{F +G}{} ein paralleles Vektorfeld.
}}
\faktzusatz {}
\faktzusatz {}

 }
{ Siehe [[Hyperfläche/Kurve/Paralleles Vektorfeld/Eigenschaften/Fakt/Beweis/Aufgabe|Aufgabe 6.3]]. }

Zu einer differenzierbaren Kurve
\maabbdisp {\gamma} { I } { Y
} {}
ist insbesondere das Geschwindigkeitsfeld $\gamma'$ ein Vektorfeld längs $\gamma$, daher ist das Konzept, ob dieses Feld parallel ist, anwendbar.

__NOEDITSECTION__

\inputfaktbeweis
{Differenzierbare Hyperfläche/Differenzierbare Kurve/Ableitung parallel/Geodätische/Fakt}
{Lemma}
{}
{

\faktsituation {Es sei

\mavergleichskette
{\vergleichskette
{ Y
}
{ \subseteq }{ W
}
{ \subseteq }{ \R^n
}
{    }{
}
{   }{
}
}
{}{}{,}
$W$
\definitionsverweis {offen}{}{,}
eine
\definitionsverweis {differenzierbare Hyperfläche}{}{}
und sei
\maabb {\gamma} { I } { Y
} {}
eine zweimal
\definitionsverweis {stetig differenzierbare Kurve}{}{.}}
\faktfolgerung {Dann ist $\gamma$ genau dann eine
\definitionsverweis {geodätische Kurve}{}{,}
wenn das Ableitungsfeld
\zusatzklammer {Geschwindigkeitsfeld} {} {} $\gamma'$ ein
\definitionsverweis {paralleles Vektorfeld}{}{}
längs $\gamma$ ist.}
\faktzusatz {}
\faktzusatz {}

}
{

Es sei
\maabbdisp {\gamma^{\prime \prime}} { I } { \R^n
} {}
die zweite Ableitung. Nach Definition ist $\gamma$ eine geodätische Kurve wenn
\mathl{\gamma^{\prime \prime} (t)}{} stets senkrecht auf dem Tangentialraum $T_{\gamma(t)} Y$ ist, was genau dann der Fall ist, wenn die orthogonale Projektion von $\gamma^{\prime \prime}(t)$ auf den Tangentialraum gleich $0$ ist. Dies ist äquivalent zu

\mavergleichskettedisp
{\vergleichskette
{ (\nabla_\gamma \gamma')(t)
}
{ =} { \pi_{\gamma(t)} \left(  \gamma^{\prime \prime} (t)  \right)
}
{ =} { 0
}
{ } {
}
{ } {
}
}
{}{}{}
für alle

\mavergleichskette
{\vergleichskette
{ t
}
{ \in }{ I
}
{  }{
}
{    }{
}
{   }{
}
}
{}{}{,}
was bedeutet, dass $\gamma'$ ein
\definitionsverweis {paralleles Vektorfeld}{}{}
längs $\gamma$ ist.

}

Die folgende Aussage ist die Grundlage für die Existenz des Paralleltransportes.
__NOEDITSECTION__

\inputfaktbeweis
{Hyperfläche/Kurve/Anfangstangentenvektor/Paralleles Vektorfeld/Existenz/Fakt}
{Satz}
{}
{

\faktsituation {Es sei

\mavergleichskette
{\vergleichskette
{ Y
}
{ \subseteq }{ W
}
{ \subseteq }{\R^n
}
{    }{
}
{   }{
}
}
{}{}{,}
$W$
\definitionsverweis {offen}{}{,}
eine
\definitionsverweis {differenzierbare Hyperfläche}{}{}
und sei
\maabb {\gamma} { I } { Y
} {}
eine
\definitionsverweis {differenzierbare Kurve}{}{.}
Es sei

\mavergleichskette
{\vergleichskette
{ t_0
}
{ \in }{ I
}
{  }{
}
{    }{
}
{   }{
}
}
{}{}{,}

\mavergleichskette
{\vergleichskette
{ \gamma(t_0)
}
{ = }{ P
}
{  }{
}
{    }{
}
{   }{
}
}
{}{}{}
und sei

\mavergleichskette
{\vergleichskette
{ v
}
{ \in }{ T_PY
}
{  }{
}
{    }{
}
{   }{
}
}
{}{}{}
ein
\definitionsverweis {Tangentialvektor}{}{.}}
\faktfolgerung {Dann gibt es ein eindeutig bestimmtes
\definitionsverweis {tangentiales Vektorfeld}{}{}
$F$ längs $\gamma$, das
\definitionsverweis {parallel}{}{}
ist und

\mavergleichskettedisp
{\vergleichskette
{ F(t_0)
}
{ =} { v
}
{ } {
}
{ } {
}
{ } {
}
}
{}{}{}
erfüllt.}
\faktzusatz {}
\faktzusatz {}

}
{

Wir suchen nach einer differenzierbaren Abbildung
\maabbdisp {F} { I } { \R^n
} {,}
das die Bedingungnen
\aufzaehlungdrei{
\mavergleichskettedisp
{\vergleichskette
{ F(t)
}
{ \in} { T_{\gamma(t)} Y
}
{ } {
}
{ } {
}
{ } {
}
}
{}{}{}
für alle

\mavergleichskette
{\vergleichskette
{ t
}
{ \in }{ I
}
{  }{
}
{    }{
}
{   }{
}
}
{}{}{,}
}{
\mavergleichskettedisp
{\vergleichskette
{ F'(t)
}
{ \in} { N_{\gamma(t)} Y
}
{ } {
}
{ } {
}
{ } {
}
}
{}{}{}
für alle

\mavergleichskette
{\vergleichskette
{ t
}
{ \in }{ I
}
{  }{
}
{    }{
}
{   }{
}
}
{}{}{,}
}{
\mavergleichskettedisp
{\vergleichskette
{ F (t_0)
}
{ =} { v
}
{ } {
}
{ } {
}
{ } {
}
}
{}{}{}
}
erfüllt. Die erste Bedingung bedeutet, dass $F(t)$ senkrecht auf dem Normaleneinheitsvektors $N (\gamma(t))$ steht, also

\mavergleichskettedisp
{\vergleichskette
{ \left\langle F( t) ,  N( \gamma(t))   \right\rangle
}
{ =} { 0
}
{ } {
}
{ } {
}
{ } {
}
}
{}{}{.}
Daher ist auch

\mavergleichskettedisp
{\vergleichskette
{ 0
}
{ =} { \left\langle F(t)  ,  N( \gamma(t))   \right\rangle'
}
{ =} { \left\langle F'(t)  ,  N( \gamma(t))   \right\rangle +  \left\langle F(t)  ,  (N( \gamma(t)))'   \right\rangle
}
{ } {
}
{ } {
}
}
{}{}{.}
Die zweite Bedingung, dass $F'(t)$ im Normalenraum liegt, bedeutet

\mavergleichskettedisp
{\vergleichskette
{ F'(t) - \left\langle F'( t) ,  N( \gamma(t))   \right\rangle N(\gamma(t) )
}
{ =} { 0
}
{ } {
}
{ } {
}
{ } {
}
}
{}{}{}
für alle

\mavergleichskette
{\vergleichskette
{ t
}
{ \in }{ I
}
{  }{
}
{    }{
}
{   }{
}
}
{}{}{,}
was wir mithilfe der vorstehenden Gleichung zu

\mavergleichskettedisp
{\vergleichskette
{ F'(t) + \left\langle F( t) ,  (N( \gamma(t)))'   \right\rangle N(\gamma(t) )
}
{ =} { 0
}
{ } {
}
{ } {
}
{ } {
}
}
{}{}{}
für alle

\mavergleichskette
{\vergleichskette
{ t
}
{ \in }{ I
}
{  }{
}
{    }{
}
{   }{
}
}
{}{}{}
bzw. zu

\mavergleichskettedisp
{\vergleichskette
{ F'(t)
}
{ =} { - \left\langle F( t) ,  (N( \gamma(t)))'   \right\rangle N(\gamma(t) )
}
{ } {
}
{ } {
}
{ } {
}
}
{}{}{}
für alle

\mavergleichskette
{\vergleichskette
{ t
}
{ \in }{ I
}
{  }{
}
{    }{
}
{   }{
}
}
{}{}{}
umformen können. Hier steht eine gewöhnliche lineare Differentialgleichung erster Ordnung für $F$, die zusammen mit der Anfangsbedingung

\mavergleichskette
{\vergleichskette
{ F(t_0)
}
{ = }{ v
}
{  }{
}
{    }{
}
{   }{
}
}
{}{}{}
nach
[[Picard Lindelöf/Lokal Lipschitz/Lokale Existenz und Eindeutigkeit/Fakt|Satz 56.2 (Analysis (Osnabrück 2021-2023))]]
eine eindeutige Lösung besitzt. Es kann also höchstens eine Lösung der Ausgangsgleichung geben. Wenn $F$ die eindeutige Lösung der Differentialgleichung ist, so liegt in der Tat eine Lösung der Ausgangsgleichung vor.

}

\inputbemerkung
{}
{

In der Situation von
[[Hyperfläche/Kurve/Anfangstangentenvektor/Paralleles Vektorfeld/Existenz/Fakt|Satz 6.9]]
ist die zu lösende Differentialgleichung das System

\mavergleichskettedisp
{\vergleichskette
{ F_i'(t)
}
{ =} { - \left(  \sum_{j = 1}^n  F_j (t) (N ( \gamma(t)))'_j  \right)   N(\gamma(t))_i
}
{ } {
}
{ } {
}
{ } {
}
}
{}{}{}
für

\mavergleichskette
{\vergleichskette
{ i
}
{ = }{ 1 , \ldots , n
}
{  }{
}
{    }{
}
{   }{
}
}
{}{}{}
oder in Matrixschreibweise

\mavergleichskettedisp
{\vergleichskette
{ \begin{pmatrix} F_1' \\ \vdots\\ F_n'   \end{pmatrix}
}
{ =} { - A \begin{pmatrix} F_1  \\ \vdots\\ F_n   \end{pmatrix}
}
{ } {
}
{ } {
}
{ } {
}
}
{}{}{}
mit den Einträgen

\mavergleichskettedisp
{\vergleichskette
{ A_{ij} (t)
}
{ =} { N(\gamma(t))_i N ( \gamma(t)))'_j
}
{ } {
}
{ } {
}
{ } {
}
}
{}{}{.}

}

\inputbeispiel{}
{

Wir betrachten den Viertelgroßkreis
\maabbeledisp {} {[0 , \pi/2] } { \R^3
} { t } { \left(  \cos  t   , \,   \sin  t  , \,   0                 \right)
} {,}
auf der Einheitskugeloberfläche $S$. Es ist

\mavergleichskette
{\vergleichskette
{ \gamma(0)
}
{ = }{ \left(  1   , \,   0  , \,   0                 \right)
}
{ = }{ P
}
{    }{
}
{   }{
}
}
{}{}{}
und

\mavergleichskette
{\vergleichskette
{ \gamma(\pi/2)
}
{ = }{ \left(  0   , \,   1  , \,   0                 \right)
}
{ = }{ Q
}
{    }{
}
{   }{
}
}
{}{}{}
mit den
\definitionsverweis {Tangentialräumen}{}{}

\mavergleichskette
{\vergleichskette
{ T_PS
}
{ = }{ \R e_2 +\R e_3
}
{  }{
}
{    }{
}
{   }{
}
}
{}{}{}
und

\mavergleichskette
{\vergleichskette
{ T_QS
}
{ = }{ \R e_1 +\R e_3
}
{  }{
}
{    }{
}
{   }{
}
}
{}{}{.}
Der nach innen zeigende Einheitsnormalenvektor längs des Weges ist

\mavergleichskette
{\vergleichskette
{ N(t)
}
{ = }{ -\gamma(t)
}
{  }{
}
{    }{
}
{   }{
}
}
{}{}{.}
Die Matrix aus
[[Hyperfläche/Kurve/Anfangstangentenvektor/Paralleles Vektorfeld/Differentialgleichung/Explizite Gestalt/Bemerkung|Bemerkung 6.10]],
die die Differentialgleichung für ein paralleles Vektorfeld beschreibt, ist

\mathdisp {\begin{pmatrix}  -\sin  t \cos  t   &   \cos  t \cos  t  &  0  \\  - \sin  t  \sin  t  &  \sin  t  \cos  t  &  0  \\ 0 &  0  &  0 \end{pmatrix}} { . }

Die konstante Funktion

\mavergleichskettedisp
{\vergleichskette
{ F(t)
}
{ =} { e_3
}
{ =} { \begin{pmatrix}  0  \\ 0\\  1   \end{pmatrix}
}
{ } {
}
{ } {
}
}
{}{}{}
ist eine Lösung. Ferner ist

\mavergleichskettedisp
{\vergleichskette
{ G(t)
}
{ =} { \begin{pmatrix}  - \sin  t  \\ \cos  t \\  0   \end{pmatrix}
}
{ } {
}
{ } {
}
{ } {
}
}
{}{}{}
eine Lösung, es ist ja einerseits

\mavergleichskettedisp
{\vergleichskette
{ G'(t)
}
{ =} { \begin{pmatrix}  - \cos  t  \\ - \sin  t \\  0   \end{pmatrix}
}
{ } {
}
{ } {
}
{ } {
}
}
{}{}{}
und andererseits

\mavergleichskettedisphandlinks
{\vergleichskettedisphandlinks
{ - \begin{pmatrix}  -\sin  t \cos  t   &   \cos  t \cos  t  &  0  \\  - \sin  t  \sin  t  &  \sin  t  \cos  t  &  0  \\ 0 &  0  &  0 \end{pmatrix} \begin{pmatrix}  - \sin  t  \\ \cos  t \\  0   \end{pmatrix}
}
{ =} { - \begin{pmatrix}  \sin^{ 2  }  t \cos  t + \cos^{ 3  }  t  \\ \sin^{ 3  }  t +  \sin  t \cos^{ 2  }  t \\  0   \end{pmatrix}
}
{ =} { - \begin{pmatrix}  \cos  t  \\ \sin  t \\  0   \end{pmatrix}
}
{ } {
}
{ } {
}
}
{}{}{.}

}

  \zwischenueberschrift{Paralleltransport}

Es sei

\mavergleichskette
{\vergleichskette
{ Y
}
{ \subseteq }{ W
}
{ \subseteq }{ \R^n
}
{    }{
}
{   }{
}
}
{}{}{,}
$W$
\definitionsverweis {offen}{}{,}
eine
\definitionsverweis {differenzierbare Hyperfläche}{}{.}
Zu Punkten

\mavergleichskette
{\vergleichskette
{ P,Q
}
{ \in }{ Y
}
{  }{
}
{    }{
}
{   }{
}
}
{}{}{}
gibt es keine natürliche Beziehung zwischen dem
\definitionsverweis {Tangentialraum}{}{}
$T_PY$ und dem Tangentialraum $T_QY$. Es sei
\maabbdisp {\gamma} {[a,b]} {Y
} {}
eine
\definitionsverweis {differenzierbare Kurve}{}{}
in $Y$ mit

\mavergleichskette
{\vergleichskette
{ \gamma(a)
}
{ = }{ P
}
{  }{
}
{    }{
}
{   }{
}
}
{}{}{}
und

\mavergleichskette
{\vergleichskette
{ \gamma(b)
}
{ = }{ Q
}
{  }{
}
{    }{
}
{   }{
}
}
{}{}{.}
Man kann sich fragen, ob es entlang dieses Weges eine sinnvolle Beziehung zwischen den Tangentialräumen gibt.

Zu einer fixierten differenzierbaren Kurve $\gamma$  in $Y$, die die Punkte
\mathkor {} {P =\gamma(a)} {und} {Q = \gamma(b)} {}
verbindet, ist durch
[[Hyperfläche/Kurve/Anfangstangentenvektor/Paralleles Vektorfeld/Existenz/Fakt|Satz 6.9]]
eine Abbildung
\maabbdisp {} {T_PY} {T_QY
} {}
festgelegt. Diese ordnet dem Anfangsvektor

\mavergleichskette
{\vergleichskette
{ v
}
{ \in }{ T_PY
}
{  }{
}
{    }{
}
{   }{
}
}
{}{}{}
den Vektor

\mavergleichskette
{\vergleichskette
{ F(b)
}
{ \in }{ T_QY
}
{  }{
}
{    }{
}
{   }{
}
}
{}{}{}
zu, wobei $F$ das eindeutig bestimmte parallele Vektorfeld längs $\gamma$ ist. Diese Abbildung heißt der \stichwort {Paralleltransport} {} längs $\gamma$. Sie wird mit
\maabbdisp {\Psi_\gamma} { T_{\gamma(a)}Y} {T_{\gamma(b)}Y
} {}
bezeichnet.

\bild{ \begin{center}
\includegraphics[width=5.5cm]{ \bildeinlesungsvg {Parallel transport sphere2} {svg} }
\end{center}
\bildtext {} }

\bildlizenz { Parallel transport sphere2.svg } {} {Silly rabbit} {en Wikipedia} {CC-by-sa 3.0} {}

__NOEDITSECTION__

\inputfaktbeweis
{Hyperfläche/Kurve/Paralleltransport/Isometrie/Fakt}
{Satz}
{}
{

\faktsituation {Es sei

\mavergleichskette
{\vergleichskette
{ Y
}
{ \subseteq }{ W
}
{ \subseteq }{ \R^n
}
{    }{
}
{   }{
}
}
{}{}{,}
$W$
\definitionsverweis {offen}{}{,}
eine
\definitionsverweis {differenzierbare Hyperfläche}{}{}
und sei
\maabb {\gamma} {[a,b]} {Y
} {}
eine
\definitionsverweis {differenzierbare Kurve}{}{}
mit

\mavergleichskette
{\vergleichskette
{ \gamma(a)
}
{ = }{ P
}
{  }{
}
{    }{
}
{   }{
}
}
{}{}{}
und

\mavergleichskette
{\vergleichskette
{ \gamma(b)
}
{ = }{ Q
}
{  }{
}
{    }{
}
{   }{
}
}
{}{}{.}}
\faktfolgerung {Dann ist der
\definitionsverweis {Paralleltransport}{}{}
längs $\gamma$ eine
\definitionsverweis {Isometrie}{}{}
\maabbdisp {} {T_PY} {T_QY
} {.}}
\faktzusatz {}
\faktzusatz {}

}
{

Zum Nachweis der Linearität seien

\mavergleichskette
{\vergleichskette
{ v,w
}
{ \in }{ T_PY
}
{  }{
}
{    }{
}
{   }{
}
}
{}{}{}
and

\mavergleichskette
{\vergleichskette
{ r,s
}
{ \in }{ \R
}
{  }{
}
{    }{
}
{   }{
}
}
{}{}{}
gegeben. Es seien
\mathkor {} {F} {bzw.} {G} {}
die gemäß
[[Hyperfläche/Kurve/Anfangstangentenvektor/Paralleles Vektorfeld/Existenz/Fakt|Satz 6.9]]
eindeutig bestimmten parallelen Vektorfelder längs $\gamma$ mit

\mavergleichskette
{\vergleichskette
{ F(a)
}
{ = }{ v
}
{  }{
}
{    }{
}
{   }{
}
}
{}{}{}
und

\mavergleichskette
{\vergleichskette
{ G(a)
}
{ = }{ w
}
{  }{
}
{    }{
}
{   }{
}
}
{}{}{.}
Nach
[[Hyperfläche/Kurve/Paralleles Vektorfeld/Eigenschaften/Fakt|Lemma 6.7]]
ist
\mathl{rF+sG}{} ein paralleles Vektorfeld mit

\mavergleichskette
{\vergleichskette
{ (rF+sG)(a)
}
{ = }{ v+ w
}
{  }{
}
{    }{
}
{   }{
}
}
{}{}{.}
Wegen der Eindeutigkeit aus
[[Hyperfläche/Kurve/Anfangstangentenvektor/Paralleles Vektorfeld/Existenz/Fakt|Satz 6.9]]
ist somit
\mathl{rF+sG}{} das parallele Vektorfeld zum Tangentialvektor
\mathl{rv+sw}{.} Daher ist

\mavergleichskettedisp
{\vergleichskette
{ \Psi_{\gamma} (rv+sw)
}
{ =} { (rF+sG)(b)
}
{ =} { r F(b) +s G(b)
}
{ =} { r \Psi_{\gamma} (v) + s \Psi_{\gamma} (w)
}
{ } {
}
}
{}{}{.}

Zum Nachweis der Verträglichkeit mit dem Skalarprodukt seien wieder

\mavergleichskette
{\vergleichskette
{ v,w
}
{ \in }{ T_PY
}
{  }{
}
{    }{
}
{   }{
}
}
{}{}{}
gegeben und es seien $F,G$ die zugehörigen parallelen Vektorfelder. Es ist

\mavergleichskettedisp
{\vergleichskette
{ \left\langle  F(t) , G(t)  \right\rangle'
}
{ =} { \left\langle  F'(t) , G(t)  \right\rangle  +  \left\langle  F(t) , G'(t)  \right\rangle
}
{ =} { 0
}
{ } {
}
{ } {
}
}
{}{}{,}
da $F,G$ tangential sind und $F',G'$ orthogonal zum Tangentialraum sind. Daher ist
\mathl{\left\langle  F(t) , G(t)  \right\rangle}{} konstant längs des Weges. Daher ist

\mavergleichskettedisp
{\vergleichskette
{ \left\langle \Psi_\gamma(v)  , \Psi_\gamma(w)  \right\rangle
}
{ =} { \left\langle  F(b)  ,  G(b)   \right\rangle
}
{ =} { \left\langle  F(a)  ,  F(a)   \right\rangle
}
{ =} { \left\langle  v  , w  \right\rangle
}
{ } {
}
}
{}{}{.}
Die Bijektivität ist damit auch klar.

}

\inputbeispiel{}
{

Wir betrachten den Viertelgroßkreis
\maabbeledisp {} {[0 , \pi/2] } { \R^3
} { t } { \left(  \cos  t   , \,   \sin  t  , \,   0                 \right)
} {,}
auf der Einheitskugeloberfläche $S$ und knüpfen an
[[Kugeloberfläche/Viertelgroßkreis/Parallele Vektorfelder/Beispiel|Beispiel 6.11]]
an. Es ist

\mavergleichskette
{\vergleichskette
{ \gamma(0)
}
{ = }{ \left(  1   , \,   0  , \,   0                 \right)
}
{ = }{ P
}
{    }{
}
{   }{
}
}
{}{}{}
und

\mavergleichskette
{\vergleichskette
{ \gamma(\pi/2)
}
{ = }{ \left(  0   , \,   1  , \,   0                 \right)
}
{ = }{ Q
}
{    }{
}
{   }{
}
}
{}{}{}
mit den
\definitionsverweis {Tangentialräumen}{}{}

\mavergleichskette
{\vergleichskette
{ T_PS
}
{ = }{ \R e_2 +\R e_3
}
{  }{
}
{    }{
}
{   }{
}
}
{}{}{}
und

\mavergleichskette
{\vergleichskette
{ T_QS
}
{ = }{ \R e_1 +\R e_3
}
{  }{
}
{    }{
}
{   }{
}
}
{}{}{.}
Nach den Berechnungen im angegebenen Beispiel gilt für den
\definitionsverweis {Paralleltransport}{}{}

\mavergleichskette
{\vergleichskette
{ \Psi_\gamma(e_3)
}
{ = }{ e_3
}
{  }{
}
{    }{
}
{   }{
}
}
{}{}{}
und

\mavergleichskette
{\vergleichskette
{ \Psi_\gamma(e_2)
}
{ = }{ - e_1
}
{  }{
}
{    }{
}
{   }{
}
}
{}{}{.}

}

Wenn
\maabbdisp {\gamma} { I } { Y
} {}
eine stetige, stückweise differenzierbare Kurve ist, so definiert man den Paralleltransport längs $\gamma$, indem man die einzelnen Paralleltransporte zu den differenzierbaren Kurvenstücken hintereinander ausführt.

 [[Kategorie:Latexseite]]
